\chapter{五次三维超曲面的镜像对称}
\SourcePage{10}
\section{五次三维超曲面、它的镜像与 COGP}
取 $\PP^4$ 中的一个五次 Calabi–Yau 超曲面，即 $V(f)\subset\PP^4$，其中 $f$ 是一个齐次五次多项式。先来计算变量 $x_0,\ldots,x_4$ 的齐次五次多项式空间的维数。这只是一个允许重复选取的组合计数问题。例如，单项式可以用竖线与分隔点编码：
\Eq{1.1}{x_0^2x_2\quad\longleftrightarrow\quad\texttt{||..|},}
\Eq{1.2}{x_0x_1^2x_2\quad\longleftrightarrow\quad\texttt{|.||.|}.}
对于五次单项式，由此得到
\Eq{1.3}{\binom95=126,}
所以
\Eq{1.4}{\dim_{\CC}\{Z(f)\subset\PP^4\}=126-1=125.}
现在再对 $\mathrm{PGL}(5)$ 取商。该群的维数为 $5\cdot5-1=24$，故
\Eq{1.5}{\dim\underbrace{\mathcal M_5}_{\text{五次超曲面的模空间}}=101.}
事实上，对于一个射影 Calabi–Yau 流形 $X$，与 $X$ 变形等价的 Calabi–Yau 流形的模空间是一个光滑的轨形，其复维数为 $h^1(\Theta_X)$，其中 $\Theta_X$ 是全纯切丛。下一讲，我们将把这个数的计算作为一项练习。对于 $V(5)\subset\PP^4$，它等于 $101$。
\Lecture{2}{2019年9月3日}
\section{五次三维超曲面}
我们将考察由齐次五次多项式的零点集 $Z(f)$ 给出的五次三维超曲面 $V(5)\subset\PP^4$。
\thmnum{theorem}{1}\begin{theorem}\label{thm:1.1}
设 $X$ 为射影 Calabi–Yau 流形，则与 $X$ 变形等价的 Calabi–Yau 流形的模空间是光滑空间，维数为 $h^1(\Theta_X)$。
\end{theorem}
我们已经用一个初等论证看出，$\dim\mathcal M_{V(5)}=101$。
\subsection{\texorpdfstring{$H^1(\Theta_X)$}{H¹(ΘX)} 的计算}
我们将在 $X=V(5)\subset\PP^4$ 的情形下计算。先从 \Term{eulerxulie}{Euler 序列}开始。我们有 $\PP^4=\operatorname{Proj}\CC[x_0,\ldots,x_4]$，而
\Eq{1.6}{x_i\partial_{x_i}=x_i\frac{\partial}{\partial x_i}}
\SourcePage{11}
是定义良好的对数向量场。Euler 序列为
\Eq{1.7}{\begin{gathered}0\longrightarrow\OO_{\PP^4}\longrightarrow\OO_{\PP^4}(1)^{\oplus5}\longrightarrow\Theta_{\PP^4}\longrightarrow0,\\1\longmapsto\sum_i e_i,\qquad e_i\longmapsto x_i\partial_{x_i}.\end{gathered}}
接下来有\Term{yufaxulie}{余法序列}
\Eq{1.8}{0\longrightarrow\II/\II^2\longrightarrow\Omega^1_{\PP^4}|_X\xrightarrow{\res_X}\Omega_X^1\longrightarrow0,}
其中 $\II$ 是 $X$ 的理想层。它的对偶序列是
\Eq{1.9}{0\longrightarrow\Theta_X\longrightarrow\Theta_{\PP^4}|_X\longrightarrow N_{X/\PP^4}\simeq\OO_X(5)\longrightarrow0.}
由于 $\II\simeq\OO(-5)$，可以算出 $\II/\II^2$。相应的限制序列为
\Eq{1.10}{0\longrightarrow\OO_{\PP^4}(-5)\longrightarrow\OO_{\PP^4}\longrightarrow\OO_X\longrightarrow0,}
\Eq{1.11}{0\longrightarrow\Theta_{\PP^4}(-5)\longrightarrow\Theta_{\PP^4}\longrightarrow\Theta_{\PP^4}|_X\longrightarrow0.}
由 \eqref{eq:1.9} 得到
\Eq{1.12}{\begin{array}{ccccccccc}H^0(\Theta_X)&\to&H^0(\Theta_{\PP^4}|_X)&\to&H^0(\OO_X(5))&\to&H^1(\Theta_X)&\to&H^1(\Theta_{\PP^4}|_X)\\0&&\CC^{24}&&\CC^{125}&&\CC^{101}&&0\end{array}.}
对于 $H^1(\Theta_{\PP^4}|_X)$，\eqref{eq:1.11} 给出
\Eq{1.13}{H^1(\Theta_{\PP^4})\longrightarrow H^1(\Theta_{\PP^4}|_X)\longrightarrow H^2(\Theta_{\PP^4}(-5)).}
Euler 序列给出
\Eq{1.14}{\underbrace{H^1(\OO_{\PP^4}(1))^{\oplus5}}_{0}\longrightarrow H^1(\Theta_{\PP^4})\longrightarrow\underbrace{H^2(\OO_{\PP^4})}_{0}.}
再将 Euler 序列与 $\OO(-5)$ 作张量积，得到所需的上同调片段
\Eq{1.15}{\underbrace{H^i(\OO_{\PP^4}(-4))^{\oplus5}}_{0}\longrightarrow H^i(\Theta_{\PP^4}(-5))\longrightarrow\underbrace{H^{i+1}(\OO_{\PP^4}(-5))}_{0}\qquad(i=0,1,2).}
\SourcePage{12}
对于 $H^0(\OO_{\PP^4}(5)|_X)$，可将 \eqref{eq:1.10} 与 $\OO(5)$ 作张量积，得到
\Eq{1.16}{\begin{array}{ccccccccc}0&\to&H^0(\OO_{\PP^4})&\to&H^0(\OO_{\PP^4}(5))&\to&H^0(\OO_X(5))&\to&H^1(\OO_{\PP^4})\\&&\CC&&\CC^{126}&&\CC^{125}&&0\end{array}.}
对于 $H^0(\Theta_{\PP^4}|_X)$，\eqref{eq:1.11} 给出
\Eq{1.17}{\underbrace{H^0(\Theta_{\PP^4}(-5))}_{0}\longrightarrow\underbrace{H^0(\Theta_{\PP^4})}_{\CC^{24}}\longrightarrow\underbrace{H^0(\Theta_{\PP^4}|_X)}_{\CC^{24}}\longrightarrow\underbrace{H^1(\Theta_{\PP^4}(-5))}_{0}.}
最后得到
\Eq{1.18}{h^1(\Theta_X)=101.}
\section{Hodge 菱形}
\subsection{Dolbeault 上同调}
设 $X$ 是紧 Kähler 流形。回忆 \Term{dolbeault}{Dolbeault 上同调} $H^{i,j}_{\bar\partial}$，它是下列复形的上同调：
\Eq{1.19}{A^{i,j}\xrightarrow{\bar\partial}A^{i,j+1}.}
其中的微分形式具有如下局部表达式：
\Eq{1.20}{\sum h_{\mu\nu}\,dz_{\mu_1}\wedge\cdots\wedge dz_{\mu_i}\wedge d\bar z_{\nu_1}\wedge\cdots\wedge d\bar z_{\nu_j},}
系数 $h_{\mu\nu}\in C^\infty$。
我们有以下事实：
\begin{enumerate}
\item Dolbeault 上同调具有典范同构
\Eq{1.21}{H^{i,j}_{\bar\partial}\simeq H^j(X,\Omega_X^i)\simeq\overline{H^{j,i}_{\bar\partial}},}
从而
\Eq{1.22}{\dim H^{i,j}_{\bar\partial}=h^{i,j}=h^{j,i}.}
\item 有\footnote{由 Serre 对偶性。}
\Eq{1.23}{\begin{aligned}H^{n-i,n-j}_{\bar\partial}&=H^{n-j}(X,\Omega_X^{n-i})\\&\simeq H^j\bigl(X,K_X\otimes(\Omega_X^{n-i})^*\bigr)^*\\&\simeq\bigl(H^{i,j}_{\bar\partial}\bigr)^* .\end{aligned}}
特别地，$h^{i,j}=h^{n-i,n-j}$。
\item 有 Hodge 分解
\Eq{1.24}{H^k(X,\CC)=H^k_{\mathrm{dR}}(X)=\bigoplus_{i+j=k}H^{i,j}_{\bar\partial}.}
\end{enumerate}
回忆 $b_k=\dim_{\CC} H^k(X,\CC)$。按照这里采用的 Calabi–Yau 定义，对于 Calabi–Yau $n$ 维流形，$b_1=0$，所以 $h^{1,0}=h^{0,1}=0$。此外，
\Eq{1.25}{H^{n,0}=H^0\bigl(X,\underbrace{K_X}_{\OO_X}\bigr)\simeq\CC.}
\SourcePage{13}
这意味着
\Eq{1.26}{h^{n,0}=h^{0,n}=1.}
若 Calabi–Yau 流形的万有覆盖 $\widetilde X\to X$ 不能分解成 Calabi–Yau 流形的非平凡乘积，我们就称它为不可约的。这里讨论的严格 Calabi–Yau 情形还满足 $H^{k,0}=0$，其中 $k=1,\ldots,n-1$。
\begin{Counterexample}{1}\label{ex:1.1}
乘积 $K3\times K3$ 不是不可约的。
\end{Counterexample}
\subsection{Hodge 菱形}
\Term{hodgelingxing}{Hodge 菱形}如下：
\Eq{1.27}{\begin{array}{ccccccccl}&&&h^{3,3}&&&&&\leftarrow H^6\\&&h^{2,3}&&h^{3,2}&&&&\leftarrow H^5\\&h^{1,3}&&h^{2,2}&&h^{3,1}&&&\vdots\\h^{0,3}&&h^{1,2}&&h^{2,1}&&h^{3,0}\\&h^{0,2}&&h^{1,1}&&h^{2,0}\\&&h^{0,1}&&h^{1,0}\\&&&h^{0,0}\end{array}}
其中真正有趣的只有中心部分：
\Eq{1.28}{\begin{array}{ccccccc}&&&1\\&&0&&0\\&0&&\color{blue}h^{2,2}&&0\\1&&\color{red}h^{1,2}&&\color{red}h^{2,1}&&1\\&0&&\color{blue}h^{1,1}&&0\\&&0&&0\\&&&1\end{array}}
这是因为
\Eq{1.29}{H^1(\Theta_X)=H^1\bigl(\Theta_X\otimes\underbrace{K_X}_{\OO_X}\bigr)=H^1(\Omega_X^{n-1})=H^{n-1,1}.}
所以问题归结为计算这些上同调群。正如我们已经看到的，在 $V(5)$ 的情形下，$h^{2,1}=101$。
\SourcePage{14}
\subsection{关于\texorpdfstring{$(1,1)$}{(1,1)}类的 Lefschetz 定理}
现在，我们在一般的紧 Kähler 流形情形下讨论。\Term{neronseveri}{Néron–Severi 群} $\mathrm{NS}(X)$ 是映射 $H^2(X,\ZZ)\to H^2(X,\CC)$ 下 $H^{1,1}$ 的原像。因此，可以直观地将其理解为 $\mathrm{NS}(X)=H^2(X,\ZZ)\cap H^{1,1}$。我们有阿贝尔层的指数序列：
\Eq{1.30}{0\longrightarrow\ZZ\xrightarrow{\cdot2\pi i}\OO_X\xrightarrow{\exp}\OO_X^\times\longrightarrow1.}
其上同调长正合序列包含
\Eq{1.31}{H^1(X,\ZZ)\longrightarrow\underbrace{H^1(X,\OO_X)}_{H^{0,1}}\longrightarrow\underbrace{H^1(X,\OO_X^\times)}_{\Pic(X)}\xrightarrow{c_1}H^2(X,\ZZ)\longrightarrow H^2(X,\OO_X).}
所以 $c_1(\Pic(X))=\mathrm{NS}(X)$。
现在我们知道了如何计算 Calabi–Yau $n$ 维流形的 $\Pic(X)$。特别地，只要 $n\ge3$，就有 $h^{1,1}=\rank_{\ZZ}\mathrm{NS}(X)$，且 $\Pic(X)\simeq\mathrm{NS}(X)$，因为在 Calabi–Yau 情形下 $\Pic^0(X)=0$。
\subsection{弱 Lefschetz 定理}
设 $X\subset\PP^n$ 为光滑超曲面，因而是 Kähler 流形。Lefschetz 定理告诉我们：
\Eq{1.32}{H^k(\PP^n,\ZZ)\longrightarrow H^k(X,\ZZ)}
当 $k<n-1=\dim X$ 时是同构，当 $k=n-1$ 时是单射。
对于这里的 Calabi–Yau 三维超曲面，$0=H^1(\PP^4,\ZZ)=H^1(X,\ZZ)$，且 $H^2(\PP^4,\ZZ)\simeq\ZZ\xrightarrow{\sim}H^2(X,\ZZ)$，所以 $\mathrm{NS}(X)\simeq\ZZ$。它由超平面类 $c_1(\OO(1))$ 生成；将其限制到 $X$ 上，就得到 $X$ 上的丰线丛，而这个线丛生成整个群。
因此，今天我们了解到，五次超曲面的 Hodge 菱形中有趣的部分是
\Eq{1.33}{\begin{array}{ccc}&1\\101&&101\\&1\end{array},}
而它的镜像五次超曲面则具有
\Eq{1.34}{\begin{array}{ccc}&101\\1&&1\\&101\end{array}.}
所以镜像具有一个很大的 Picard 群。
\Lecture{3}{2019年9月5日}
\section{Lefschetz 超平面定理}
这是对上一讲的补充。设 $X\subset\PP^n$ 是光滑超曲面，则
\Eq{1.35}{H^k(\PP^n,\ZZ)\xrightarrow{\res}H^k(X,\ZZ)}
当 $k<n-1$ 时是同构，当 $k=n-1$ 时是单射。
\SourcePage{15}
\section{镜像五次超曲面}
我们应当得到彼此镜像的 Hodge 数。回忆，原五次超曲面的 Hodge 菱形中有趣的部分是
\Eq{1.36}{\begin{array}{ccc}&1\\101&&101\\&1\end{array},}
而镜像的菱形应当是
\Eq{1.37}{\begin{array}{ccc}&101\\1&&1\\&101\end{array}.}
所以 $h^{2,1}(Y)=1$，即它具有一维模空间；但 $h^{1,1}(Y)=101$，因此 $\Pic(Y)$ 很大。
\subsection{构造}
这一构造的物理动机，是对与 Fermat 五次超曲面有关的“极小共形场论”作轨形商。Fermat 五次超曲面的方程为 $x_0^5+\cdots+x_4^5=0$。群 $(\ZZ/5)^5$ 在 $\PP^4$ 上对角作用。因为其中的对角子群 $\ZZ/5$ 作用平凡，所以得到 $(\ZZ/5)^4=(\ZZ/5)^5/(\ZZ/5)$ 的有效作用。现在对其中的一个子群 $G=(\ZZ/5)^3\subset(\ZZ/5)^4$（即固定一个坐标、例如令 $\zeta_0=1$ 的那些元素）取有限商
\Eq{1.38}{\bar Y=X/(\ZZ/5)^3.}
该作用不是自由的，所以商不是流形，而是带有某些轨形奇点的空间。$X$ 本身则由 Jacobi 判别准则可知是光滑的。某些点的稳定子 $G_y\subset(\ZZ/5)^3$ 非平凡。分两种情形：
\begin{itemize}
\item $x_i=x_j=0$，其中 $i\ne j$。此时 $G_y\simeq\ZZ/5$。这给出五次曲线 $\widetilde C_{ij}=Z(x_i,x_j)\subset X$。局部作用为
$\zeta(z_1,z_2,z_3)=(\zeta z_1,\zeta^{-1}z_2,z_3)$，因此横向出现 $uv=w^5$ 所描述的 $A_4$ 奇点。
\item $x_i=x_j=x_k=0$，其中 $i,j,k$ 两两不同。此时 $G_y\simeq(\ZZ/5)^2$，并有
\Eq{1.39}{\widetilde P_{i,j,k}\longrightarrow P_{ijk}\in\bar Y.}
局部作用为
\Eq{1.40}{(\zeta,\xi)\cdot(z_1,z_2,z_3)=(\zeta\xi z_1,\zeta^{-1}z_2,\xi^{-1}z_3).}
\end{itemize}
\begin{Note}{1.1}在 $\bar Y$ 中，
\Eq{1.41}{C_{01}=Z(x_0,x_1,x_2^5+x_3^5+x_4^5)/(\ZZ/5)^3\simeq Z(u+v+w)\simeq\PP^1\subset\PP^2.}
每一条 $C_{ij}$ 都同构于 $\PP^1$。
\end{Note}
我们希望在局部通过吹起消解这些奇点。但若还希望留在射影代数簇的范畴中，这就很微妙：吹起以后，未必仍有丰线丛。因此还必须证明一个额外的结论。
\thmnum{proposition}{2}\begin{proposition}\label{pro:1.2}存在一个射影消解 $Y\to\bar Y$。\end{proposition}
\SourcePage{16}
采用环面方法最有效，但也可以直接完成这一构造。
来数一数 $Y$ 中相互独立的例外除子。每条 $C_{ij}$ 上方有 $4$ 个，每个 $P_{ijk}$ 上方有 $6$ 个。因此，来自 $C_{ij}$ 的有 $40$ 个，来自 $P_{ijk}$ 的有 $60$ 个，总共 $100$ 个。它们与超平面类一起生成 $H^2$。
\thmnum{proposition}{3}\begin{proposition}\label{pro:1.3}$h^{1,1}(Y)=101$，$h^{2,1}(Y)=1$。\end{proposition}
S. S. Roan 给出了直接证明，Batyrev 则利用环面方法证明了这一结论。
\subsection{镜像族}
上述对 $G$ 取商的构造，可推广到所谓的 \Term{dworkzu}{Dwork 族}。具体而言，令 $X_\psi=V(f_\psi)$，其中
\Eq{1.42}{f_\psi=x_0^5+\cdots+x_4^5-5\psi x_0x_1\cdots x_4,}
$\psi$ 是复参数。于是有
\Eq{1.43}{\begin{tikzcd}(\ZZ/5)^4=\{(\zeta_0,\ldots,\zeta_4)\mid\prod\zeta_i=1\}\arrow[r,hook]\arrow[d,two heads]\&(\ZZ/5)^5\arrow[d,two heads]\\G=(\ZZ/5)^3\arrow[r,hook]\&(\ZZ/5)^4\end{tikzcd}}
以及
\Eq{1.44}{\begin{tikzcd}X_\psi=Z(f_\psi)\arrow[d,"{/G}"]\\\bar Y_\psi\&Y_\psi\arrow[l]\end{tikzcd}.}
注意，当 $\zeta^5=1$ 时，$Y_\psi\simeq Y_{\zeta\psi}$。因此有
\Eq{1.45}{\begin{tikzcd}\mathcal Y'\arrow[r,"{/(\ZZ/5)}"]\arrow[d]\&\mathcal Y\arrow[d]\\\PP^1_\psi\arrow[r,"{/(\ZZ/5)}"]\&\PP^1_z\end{tikzcd},}
其中 $z=(5\psi)^{-5}$。这是一个 Calabi–Yau 三维流形族，在 $z\ne0,5^{-5},\infty$ 时其纤维光滑。特殊纤维如下：
\begin{itemize}
\item $z=0$：方程退化为 $x_0\cdots x_4=0$，即
\Eq{1.46}{X_{z=0}=\bigcup_{i=0}^4\PP^3_i\subset\PP^4,}
是五个坐标超平面的并。
\item $z=5^{-5}$，即 $\psi=1$。这对应一个三维 $A_1$ 奇点。$X_1$ 有 $125$ 个三维 $A_1$ 奇点，局部方程为 $x^2+y^2+z^2+w^2=0$。这些奇点位于同一个 $(\ZZ/5)^3$ 轨道中。因此 $Y_{5^{-5}}$ 有一个三维 $A_1$ 奇点，有时称为“conifold”。
\item $z=\infty$，即 $\psi=0$，对应 Fermat 五次超曲面。此时不再要求作用中的单位根满足乘积为 $1$，因而出现额外的 $\ZZ/5$ 对称性，所以它实际上是一个轨形点。
\end{itemize}
现在，至少从物理学的角度看，我们已经完成了任务。\footnote{Siebert 教授说，数学家或许也该如此：不要担忧太多，且看生活会带来什么。}
\SourcePage{17}
\section{Yukawa 耦合}
\subsection{A 模型}
A 模型是辛几何一侧的模型，我们在这一侧研究五次超曲面。令 $h=\mathrm{PD}(\text{超平面})$，则 $H^2(X,\ZZ)=\ZZ\cdot h$。Yukawa 耦合为
\Eq{1.47}{\langle h,h,h\rangle_A=5+\sum_{d>0}N_d d^3q^d.}
在当时的历史背景下，$N_d$ 究竟是什么还不清楚，因为 Gromov–Witten 理论正几乎同时发展起来。如今我们知道，$N_d$ 是次数为 $d$ 的有理曲线（即亏格 $g=0$ 的曲线）的 Gromov–Witten 计数。如果用 $n_d$ 表示本原计数，则同一个表达式可以写成
\Eq{1.48}{5+\sum_{d>0}d^3n_d\underbrace{\frac{q^d}{1-q^d}}_{\text{次数为 }d\text{ 的曲线的多重覆盖}}\in\QQ[[q]].}
\subsection{B 模型}
现在考虑镜像五次超曲面 $Y_z$，其中 $z=(5\psi)^{-5}$。此时 Yukawa 耦合定义为
\Eq{1.49}{\langle\partial_z,\partial_z,\partial_z\rangle_B:=\int_{Y_z}\Omega^\nu(z)\wedge\partial_z^3\Omega^\nu(z),}
$\Omega^\nu$ 是“归一化”的全纯体积形式，即
\Eq{1.50}{\int_{\beta_0}\Omega^\nu=\text{常数},}
其中 $\beta_0\in H_3(Y,\ZZ)$。
我们需要在这一侧找到与辛结构模空间上的向量场 $h$ 对应的某种镜像。真正要使用的是指数化的量 $e^{2\pi i h}$。结果表明，在镜像对称的这一侧，它对应于 $\partial/\partial w$，也就是 $Y$ 的复结构模空间上的一个向量场。其中
\Eq{1.51}{w=\int_{\beta_1}\Omega^\nu(z),}
$\beta_1\in H_3(Y,\ZZ)$。
\subsection{镜像对称}
现在，镜像对称的实际陈述为
\Eq{1.52}{\langle h,h,h\rangle_A=\langle\partial_w,\partial_w,\partial_w\rangle_B,}
这里 $q=e^{2\pi i w(z)}$，且 $q=c\,z+O(z^2)$。于是，我们要做以下事情：
\begin{enumerate}
\item 写出全纯三形式，这还不算太难。
\item 计算归一化所用的周期积分，这也还不算太难。
\item 计算第二个积分，这就更困难了。
\item 计算 Yukawa 耦合。
\end{enumerate}
\SourcePage{18}
\subsection{周期的计算}
Mark Gross 的 Nordfjordeid 讲义中有关于这一计算的说明。回忆，我们有
\Eq{1.53}{\begin{tikzcd}\&Y_\psi\arrow[d]\\X_\psi\arrow[r,"{/G}"]\&\bar Y_\psi\end{tikzcd}.}
已知 $H_3(Y_\psi,\ZZ)\simeq\ZZ^4$。在 $\psi=\infty$ 附近，也就是\Term{dafujiegoujixian}{大复结构极限}附近，有一个\Term{xiaoshixunhuan}{消失循环}。
这与 Siebert 教授最喜欢的退化图景相似。考虑 $zw=t$：当 $t=0$ 时，它看上去是两个圆盘在一点相交；当 $t\ne0$ 时，则看上去是一个圆柱。考察 Dehn 扭转，就可以看到一个 $S^1$ 在这一退化过程中被压缩为一点。在高维中也有类似的情形。
具体来说，我们的消失循环是 $\beta_0=T^3$。局部有
\Eq{1.54}{u_1\cdots u_4=z,\qquad\beta_0=\{\,|u_1|=\cdots=|u_4|=|z|^{1/4},\ \arg(u_1\cdots u_4)=\arg z\,\},}
其中 $u_i$ 是全纯坐标。
将其提升到 $X_\psi$，就得到一个显式的三维环面
\Eq{1.55}{\begin{aligned}T=\{\,&|x_0|=|x_1|=|x_2|=\delta\ll1,\\&x_3=x_3(x_0,x_1,x_2)\text{ 是 }f_\psi(x_0,x_1,x_2,x_3,1)=0\text{ 的解},\\&\text{该近零解唯一，且当 }z\to0\text{ 时 }x_3\to0\,\}.\end{aligned}}
\Lecture{4}{2019年9月10日}
回忆我们有
\Eq{1.56}{\begin{tikzcd}\&Y_\psi\arrow[d]\\X_\psi\arrow[r,"{/G}"]\&\bar Y_\psi\end{tikzcd},}
其中 $G=(\ZZ/5\ZZ)^3$。上一讲讨论了：
\begin{enumerate}
\item 消失循环 $T^3$，以及 $\beta_0\in H_3(Y_\psi,\ZZ)$；
\item 全纯三形式；
\item 归一化；
\item 其他周期；
\item 典范坐标与镜像映射。
\end{enumerate}
\section{全纯三形式}
我们将把全纯三形式构造成 $\PP^4$ 上一个亚纯或有理四形式的留数；该四形式沿 $X_\psi$ 有极点：
\Eq{1.57}{\Omega(\psi)=5\psi\operatorname{Res}_{X_\psi}\frac{\widetilde\Omega}{f_\psi}\in\Gamma(X_\psi,\Omega^3_{X_\psi}),}
其中
\Eq{1.58}{\widetilde\Omega=\sum_{i=0}^4(-1)^i x_i\,dx_0\wedge\cdots\wedge\widehat{dx_i}\wedge\cdots\wedge dx_4.}
在局部坐标 $x_4=1$、$\partial_{x_3}f_\psi\ne0$ 下，
\Eq{1.59}{\Omega(\psi)=\left.5\psi\frac{dx_0\wedge dx_1\wedge dx_2}{\partial_{x_3}f_\psi}\right|_{X_\psi}.}
\SourcePage{19}
\section{归一化}
现在处理归一化。令 $\varphi_0=\int_{\beta_0}\Omega(\psi)$，则
\Eq{1.60}{\widetilde\Omega=\varphi_0^{-1}\Omega(\psi)}
就是归一化后的形式。利用留数，可以计算
\Eq{1.61}{2\pi i\int_{\beta_0}\Omega(\psi)=\int_{T^4}\frac{5\psi\,dx_0\,dx_1\,dx_2\,dx_3}{f_\psi},}
\Eq{1.62}{=\int_{T^4}\frac{dx_0\cdots dx_3}{x_0\cdots x_3}\frac{1}{\dfrac{1+x_0^5+\cdots+x_3^5}{5\psi x_0\cdots x_3}-1},}
\Eq{1.63}{=-\sum_{n\ge0}\int_{T^4}\frac{dx_0\cdots dx_3}{x_0\cdots x_3}\frac{(1+x_0^5+\cdots+x_3^5)^n}{(5\psi)^n(x_0\cdots x_3)^n},}
\Eq{1.64}{=-\sum_{n\ge0}\int_{T^4}\frac{dx_0\cdots dx_3}{x_0\cdots x_3}\frac{(1+x_0^5+\cdots+x_3^5)^{5n}}{(5\psi)^{5n}(x_0\cdots x_3)^{5n}}.}
要对积分有贡献，分子中每个变量的幂次必须与分母的幂次相同，因此只有原求和指标为 $5$ 的倍数的项能够留下。通过组合计数，得到
\Eq{1.66}{2\pi i\int_{\beta_0}\Omega(\psi)=-(2\pi i)^4\sum_{n\ge0}\frac{(5n)!}{(n!)^5(5\psi)^{5n}}=: \varphi_0(z),}
其中 $z=(1/(5\psi))^5$。
系数 $(5n)!/(n!)^5$，就是单项式
\Eq{1.67}{x_0^{5n}\cdots x_3^{5n}\quad\text{在}\quad(1+x_0^5+\cdots+x_3^5)^{5n}}
展开式中出现的次数。
\section{其他周期}
有一种方法称为 Griffiths 的极点阶数约化，它涉及 \Term{picardfuchs}{Picard–Fuchs 方程}。
$H^3(Y_\psi,\CC)$ 在局部是一个秩为 $4$ 的常值系统，因此给出平凡的全纯向量丛
\Eq{1.68}{\begin{tikzcd}E=U\times\CC^4\arrow[d]\\U\arrow[r,hook]\&\mathcal M_Y\end{tikzcd}.}
它带有一个平坦联络 $\nabla^{\mathrm{GM}}$，称为 \Term{gaussmanin}{Gauss–Manin 联络}。在每一点上，
\Eq{1.69}{E_\psi=\bigoplus_{p+q=3}H^{p,q}(Y_\psi),}
并且已经看到 $\Omega(\psi)\in H^{3,0}$。现在考虑
\Eq{1.70}{\Omega(z),\quad\partial_z\Omega(z),\quad\partial_z^2\Omega(z),\quad\partial_z^3\Omega(z),\quad\partial_z^4\Omega(z).}
它们由一个系数全纯的四阶常微分方程联系起来，这就是 Picard–Fuchs 方程。
\SourcePage{20}
\subsection{方程的推导}
取 $X=Z(f_\psi)\subset\PP^4$。可以利用具有更高阶极点的形式，产生更多的三形式。考虑长正合序列
\Eq{1.71}{\begin{gathered}H^4(\PP^4,\CC)\longrightarrow H^4(\PP^4\setminus X,\CC)\longrightarrow H^5(\PP^4,\PP^4\setminus X;\CC)\longrightarrow H^5(\PP^4,\CC),\\H^5(\PP^4,\PP^4\setminus X;\CC)\overset{\text{切除}}\simeq H^5(U,U\setminus X;\CC)\simeq H^5(U,\partial U;\CC)\\\overset{\mathrm{LD}}\simeq H_3(U\setminus\partial U;\CC)\simeq H_3(X,\CC)\simeq H^3(X,\CC).\end{gathered}}
其中 $X\subset U$，$U$ 是管状邻域。这里使用的 Lefschetz 对偶性是 $H^q(M,\partial M)=H_{n-q}(M\setminus\partial M)$，最后一步则使用 Poincaré 对偶性。于是，从高阶极点的形式出发，可以得到某个 $H^3$ 中的类。在我们的情形下，再对其求导；对于某些类，我们知道它们应当为零，这便给出了若干方程。
\subsection{Griffiths 的极点阶数约化}
如果有
\Eq{1.72}{\frac{g\widetilde\Omega}{f^l}\in H^0(\PP^4\setminus X,\Omega^4_{\PP^4\setminus X}),}
则必须满足 $\deg g=5l-5$。前面使用的是 $l=1$，因而分子 $g$ 是常数。恰当形式具有如下表达式：
\Eq{1.73}{d\left(\frac1{f^l}\sum_{i<j}(-1)^{i+j}(x_ig_j-x_jg_i)\,dx_0\wedge\cdots\wedge\widehat{dx_i}\wedge\cdots\wedge\widehat{dx_j}\wedge\cdots\wedge dx_4\right)}
\Eq{1.74}{=\left(l\sum_jg_j\partial_{x_j}f-f\sum_j\partial_{x_j}g_j\right)\frac{\widetilde\Omega}{f^{l+1}}.}
若分子属于 Jacobi 理想 $J(f)=(\partial_{x_i}f)$，就能把它表示成 $l\sum_jg_j\partial_{x_j}f$；模去一个恰当形式后，留下的项含有一个因子 $f$，从而可以降低极点阶数。也就是说，第一项在除以 $f^{l+1}$ 后有 $l+1$ 阶极点，而第二项的极点阶数只有 $l$。关键是：分子 $g\in J(f)$ 时，可以降低 $l$。
因此，算法如下：计算 $\Omega(z),\partial_z\Omega(z),\partial_z^2\Omega(z),\ldots,\partial_z^4\Omega(z)$。最后一项可以写成 $g\widetilde\Omega/f_\psi^5$，其中 $g\in J(f_\psi)$。再利用 \eqref{eq:1.74} 的关系，将它表示为较低阶 $\partial_z^i\Omega(z)$ 的线性组合。
\thmnum{proposition}{4}\begin{proposition}\label{pro:1.4}
任一周期
\Eq{1.75}{\varphi=\int_\alpha\Omega(\psi)}
都满足常微分方程
\Eq{1.76}{\bigl[\theta^4-5z(5\theta+1)(5\theta+2)(5\theta+3)(5\theta+4)\bigr]\varphi(z)=0,}
其中 $\theta=z\partial_z$。
\end{proposition}
\SourcePage{21}
\begin{Note}{1.2}对于
\Eq{1.77}{\varphi=\varphi_0=\sum_{n\ge0}\frac{(5n)!}{(n!)^5}z^n,}
很容易直接验证 \eqref{eq:1.76}。
\end{Note}
\eqref{eq:1.76} 在 $z=0$ 处是具有正则奇点的常微分方程。写成一阶系统，就是
\Eq{1.78}{\theta\varphi(z)=A(z)\varphi(z),}
其中 $\varphi(z)\in\CC^s$。
\thmnum{theorem}{5}\begin{theorem}\label{thm:1.5}
\eqref{eq:1.78} 有一个如下形式的基本解矩阵：
\Eq{1.79}{\Phi(z)=S(z)z^R,}
其中 $S(z)\in M(s,\OO_0)$、$R\in M(s,\CC)$，且
\Eq{1.80}{z^R=I+(\log z)R+\frac{(\log z)^2}{2!}R^2+\cdots.}
如果特征值之差不是非零整数，则可取 $R=A(0)$。
\end{theorem}
对于 \eqref{eq:1.76}，
\Eq{1.81}{A(0)\simeq\begin{pmatrix}0&1&0&0\\0&0&1&0\\0&0&0&1\\0&0&0&0\end{pmatrix},}
而 $S=(\psi_0,\ldots,\psi_3)$，其中 $\psi_i\in\OO_{\CC,0}^4$。由此得到一组基本解：
\Eq{1.82}{\varphi_0(z)=\psi_0(z)\qquad\text{（单值）},}
\Eq{1.83}{\varphi_1(z)=\psi_0(z)\log z+\psi_1(z),}
\Eq{1.84}{\varphi_2(z)=\tfrac12\psi_0(z)(\log z)^2+\psi_1(z)\log z+\psi_2(z),}
\Eq{1.85}{\varphi_3(z)=\tfrac16\psi_0(z)(\log z)^3+\cdots+\psi_3(z).}
这与\Term{danzhi}{单值变换}有关。特别地，$z^{A(0)}$ 的单值变换反映了 $H^3(Y_z,\CC)$ 绕 $z=0$，即 $\psi=\infty$ 的单值变换算子 $T$。事实上，可以选取 $H_3(Y_z,\QQ)$ 的一个辛基 $\beta_0,\beta_1,\alpha_1,\alpha_0$。记 $N=T-I$，则有
\Eq{1.86}{\alpha_0\longmapsto\alpha_1\longmapsto\beta_1\longmapsto\beta_0\longmapsto0.}
相应的周期为
\Eq{1.87}{\varphi_0=\int_{\beta_0}\Omega(z),\quad\varphi_1=\int_{\beta_1}\Omega(z),\quad\varphi_2=\int_{\alpha_1}\Omega(z),\quad\varphi_3=\int_{\alpha_0}\Omega(z).}
\section{典范坐标与镜像映射}
从解的形式来看，我们没有太多选择。所选的解在指数化后应当表现得像 $z$。事实上，\Term{jingxiangyingshe}{镜像映射}给出的典范坐标为
\Eq{1.88}{q=e^{2\pi i w},}
其中
\Eq{1.89}{w=\frac{\int_{\beta_1}\Omega(z)}{\int_{\beta_0}\Omega(z)}=\int_{\beta_1}\widetilde\Omega(z).}
\SourcePage{22}
于是 $\varphi_1(z)=\varphi_0(z)\log z+\psi_1(z)$。很容易通过 \eqref{eq:1.76} 的级数解求出
\Eq{1.90}{\psi_1(z)=5\sum_{n\ge1}\frac{(5n)!}{(n!)^5}\left(\sum_{j=n+1}^{5n}\frac1j\right)z^n.}
其中还可相差一个常数 $c_2$。
\Lecture{5}{2019年9月12日}
上一讲，我们学会了怎样计算这些周期、得到典范坐标，以及利用 Picard–Fuchs 方程降低极点阶数。剩下的主题是 Yukawa 耦合。
\section{Yukawa 耦合}
我们要计算
\Eq{1.91}{\langle\partial_z,\partial_z,\partial_z\rangle_B=\int_{Y_z}\widetilde\Omega(z)\wedge\partial_z^3\widetilde\Omega(z),}
其中 $\widetilde\Omega(z)=\Omega(z)/\varphi_0(z)$。
引入辅助量
\Eq{1.92}{W_k=\int_{Y_z}\Omega(z)\wedge\partial_z^k\Omega(z),\qquad k=0,\ldots,4.}
所以，我们真正需要的只是 $W_3$。把 Picard–Fuchs 方程改写为
\Eq{1.93}{\left(\frac{d^4}{dz^4}+\sum_{k=0}^3c_k\frac{d^k}{dz^k}\right)\Omega(z)=0,}
得到
\Eq{1.94}{W_4+\sum_{k=0}^3c_kW_k=0.}
\subsection{Griffiths 横截性}
现在需要引入一项重要信息，称为 \Term{griffithshengjie}{Griffiths 横截性}。它与 Hodge 结构变动的定义有关。设 $U$ 是模空间中的开集，定义递减滤过
\Eq{1.95}{F=H^3(Y_z,\ZZ)\otimes_{\ZZ}\OO_U=F^0\supset F^1\supset F^2\supset F^3,}
其中 $F$ 的秩为 $4$，且逐点有
\Eq{1.96}{F_z^k=\bigoplus_{p\ge k}H^{p,3-p}(Y_z).}
这就是 Hodge 滤过。注意，$\Omega\in\Gamma(U,\pi_*\Omega^3_{\mathcal Y/U})$，而 $\pi_*\Omega^3_{\mathcal Y/U}=F^3\simeq\OO_U$。既然逐点看上去有直和分解，为什么仍采用滤过的写法？从抽象层面看，我们有
\Eq{1.97}{\nabla^{\mathrm{GM}}F^k\subseteq F^{k-1}\otimes\Omega^1_{\mathcal M}.}
这一包含关系来自 $\nabla^{\mathrm{GM}}$ 的定义和构造，以及 Hodge 与 Dolbeault 理论。此外，除非 $p+p'=3=q+q'$，否则 $H^{p,q}\perp H^{p',q'}$；这可以由积分 $\int_Y\alpha\wedge\beta=0$ 看出。
\SourcePage{23}
这些事实合起来给出 $W_0=W_1=W_2=0$。特别地，
\Eq{1.98}{0=\frac{d^2W_2}{dz^2}=\cdots=2W_3'-W_4.}
由 \eqref{eq:1.94} 得知
\Eq{1.99}{W_3'+\frac12c_3W_3=0.}
现在计算出
\Eq{1.100}{c_3(z)=\frac6z-\frac{2\cdot5^5}{1-5^5z}.}
分离变量，得到
\Eq{1.101}{W_3=\frac{c_1}{(2\pi i)^3z^3(5^5z-1)},}
其中 $c_1$ 是积分常数。最后改用 $q=e^{2\pi i w}$，这里 $w=\varphi_1(z)/\varphi_0(z)$。展开后得到
\Eq{1.102}{\begin{aligned}\langle\partial_w,\partial_w,\partial_w\rangle_B={}&-c_1-575\frac{c_1}{c_2}q-\frac{1950750}{2}\frac{c_1}{c_2^2}q^2\\&-\frac{10277490000}{6}\frac{c_1}{c_2^3}q^3+\cdots.\end{aligned}}
它应当等于
\Eq{1.103}{5+\sum_{d\ge1}d^3n_d\frac{q^d}{1-q^d},}
即
\Eq{1.104}{5+n_1q+(8n_2+n_1)q^2+(27n_3+n_1)q^3+\cdots.}
这预测 $n_1=2875$，该数早在经典研究中就已知；$n_2=609250$ 也正确，这在上述发展前几年，即 1986 年，已由 Katz 证明。接着得到 $n_3=317206375$。这个数最初与直接计算的结果不一致，但后来发现直接计算中有一个错误，修正后也一致了。Givental 于 1996 年给出了第一个证明。
由此得到的启示如下。\footnote{此外，我们还了解到：计算确实很困难……}
\begin{enumerate}
\item 这一预测依赖于大复结构极限，因为它与 Kähler 锥有关：
\Eq{1.105}{\mathcal K_X=\{\,[\omega]\in H^2_{\mathrm{dR}}(X)\mid\omega\text{ 是 Kähler 形式}\,\}.}
特别地，$H^3(Y)=\bigoplus_{p=0}^3H^{p,3-p}$ 中的单值变换，对应于 $\bigoplus_{p=0}^3H^{p,p}(X)$ 上与 $[\omega_X]$ 的杯积。不过请注意，$\langle\,\rangle_B$ 是在整个 $\mathcal M_B$ 上定义的！
\item 通过轨形商构造镜像，是五次超曲面的特殊情形。Batyrev 与 Borisov 考察了环面簇中完全交的镜像。
\end{enumerate}
